% SPDX-License-Identifier: Apache-2.0
\section{Writing Memory in Bursts}
\label{sec:x-dmaw}

\code{LineStore} is \code{LineFetch} the other way round, and the
half an Ethernet receiver wants: words arrive on a channel and are
written to memory in bursts, rather than the core taking each one and
storing it.

Two of its differences from the fetcher are the bus rather than
choices made here.

A write burst is finished by its response and not by its last beat.
The host answers a read on the beat channel and a write on
\code{done}, so the identifier goes back when the response arrives,
which may be well after the last beat has gone out.

AXI4 puts no identifier on the beats, so a burst's beats belong to
the oldest address phase outstanding. A second burst issued before
the first one's beats were out would take them. \code{LineFetch}
keeps one burst in flight because a second buys nothing; here it is
not an optimisation being declined but a rule being obeyed.

\subsection{A frame is a count of bytes}

A line of pixels is a whole number of words. A frame is not. An
Ethernet frame is anything from 64 to 1518 bytes, so its last word
usually holds one, two or three real bytes, and writing the other
three puts whatever the channel happened to be carrying into memory
past the end of the frame.

So \code{LineStore} is told a count of bytes, not of words. It asks
for \code{(bytes + 3) / 4} beats and strobes the last one down to
the bytes that are real. Nothing about the burst changes; the beat
goes out either way, and the strobe says which lanes of it are meant.

\lstinputlisting[rangebeginprefix=//\ begin\{,rangebeginsuffix=\},%
rangeendprefix=//\ end\{,rangeendsuffix=\},includerangemarker=false,%
linerange=store-store,caption={\code{lib/parts/src/dma.rs}: what it
remembers, including how many bytes of the last word are
real.},label={lst:x-dmaw}]{lib/parts/src/dma.rs}

\subsection{Checking a write by reading it back}

The run writes sixty-four words, of which the last holds one byte,
and then reads the memory and checks it. Reading it back is the
point: a write path watched only on the bus can be wrong in ways that
look perfectly well formed going past.

Three things are asserted. Every word is the one that belongs at its
address, since each word written is its own index. Nothing past the
run was touched, which is how a burst that overran its count would
show. And the last word holds one byte of the frame over three bytes
that were there before.

That last check only works because the memory is poisoned first. The
word is filled with \code{0xaaaaaa00} before the run, so a strobe
that wrote the whole word wipes it and a strobe that wrote one byte
leaves \code{0xaaaaaa3f}. Without the poison both cases leave zeros
in those three bytes and look identical, and the check would pass on
an engine that ignored the strobe entirely.

\subsection{Where a test model goes to be wrong}

The memory model in the example had to be corrected twice before it
could see what it was for, and both faults were the same kind: a
model permissive enough that a broken engine would have passed.

It answered one beat per request at first, which makes a burst look
like a word. The words all arrive, the cycle count passes, and
nothing about bursting is demonstrated.

Then it wrote whole words and ignored the strobe, which makes a
byte-accurate engine and one that has never heard of byte enables
indistinguishable.

Neither showed up as a failure. Both showed up as a test that passed
and proved less than it claimed, which is the failure mode worth
looking for in the model rather than in the design.

\lstinputlisting[style=ascii,caption={What \code{ex\_dmaw} printed
when the build ran it.},label={lst:x-dmaw-out}]{out_dmaw.txt}
